\documentclass[11pt,letterpaper]{article}
\usepackage[top=0.7in,bottom=0.7in,left=0.75in,right=0.75in]{geometry}
\usepackage[T1]{fontenc}
\usepackage[utf8]{inputenc}
\usepackage{lmodern}
\usepackage{microtype}
\usepackage{enumitem}
\usepackage[hidelinks]{hyperref}
\usepackage{titlesec}
\setlength{\parindent}{0pt}
\setlength{\parskip}{0pt}
\pagestyle{empty}
\titleformat{\section}{\normalfont\bfseries}{}{0em}{\uppercase}[\vspace{-4pt}\hrule\vspace{4pt}]
\titlespacing{\section}{0pt}{8pt}{5pt}
\setlist[itemize]{leftmargin=1.2em,label=--,itemsep=0pt,topsep=2pt,parsep=0pt}
\begin{document}
\begin{center}
{\fontsize{16}{19}\selectfont\textbf{\MakeUppercase{Diego Castillo Pineda}}}\\[4pt]
{\small La Florida, Santiago, Chile \quad|\quad +56 9 9416 2547 \quad|\quad \href{mailto:diego.castillop11@gmail.com}{diego.castillop11@gmail.com}\\
\href{https://linkedin.com/in/diego-castillo-pineda}{linkedin.com/in/diego-castillo-pineda} \quad|\quad \href{https://github.com/diegocastillop11}{github.com/diegocastillop11}}
\end{center}

\section{Resumen Profesional}
Analista Data Science con sólida trayectoria en administración de bases de datos SQL Server, Python y automatización de procesos analíticos en entornos de alto volumen transaccional. Especializado en optimización de consultas complejas, modelado de datos y desarrollo de soluciones con Azure SQL y Machine Learning. Reduje en 40\% los tiempos de generación de reportería analítica mediante automatización con Python y SQL dinámico en Punto Ticket. Experiencia aplicada en análisis exploratorio de datos, visualización con Power BI y desarrollo de pipelines ETL para inteligencia de negocios en retail y entretenimiento masivo.

\section{Experiencia Profesional}

\textbf{Punto Ticket} \hfill Santiago, Chile\\
\textit{Analista de Base de Datos Senior} \hfill Nov. 2019 – Feb. 2026
\begin{itemize}
  \item Diseñé y optimicé stored procedures en SQL Server para procesar más de 50,000 transacciones por hora en eventos masivos, reduciendo tiempos de respuesta en 35\% mediante análisis de planes de ejecución y ajuste de índices
  \item Automaticé procesos de extracción y transformación de datos (ETL) con Python y T-SQL, disminuyendo errores manuales en 60\% y acelerando la disponibilidad de datos para análisis de comportamiento de clientes
  \item Desarrollé consultas dinámicas con PIVOT y funciones analíticas avanzadas para reportería ejecutiva, generando dashboards en Power BI que permitieron identificar patrones de compra y optimizar estrategias de pricing
  \item Integré datos de múltiples fuentes utilizando JSON, XML y APIs REST, consolidando información transaccional para modelos predictivos de demanda en Python (scikit-learn, pandas)
  \item Participé en migración de bases de datos on-premise a Azure SQL Database, implementando estrategias de backup y alta disponibilidad que aseguraron 99.9\% de uptime en períodos críticos
  \item Realicé análisis exploratorio de datos (EDA) sobre datasets de +5M registros, identificando outliers y tendencias que informaron decisiones estratégicas del área comercial
\end{itemize}
\vspace{2pt}
\textbf{Imperial S.A.} \hfill Santiago, Chile\\
\textit{Analista Funcional} \hfill Jun. 2017 – Nov. 2019
\begin{itemize}
  \item Ejecuté análisis funcional de requerimientos de negocio y traduje especificaciones a consultas SQL complejas en Oracle y SQL Server, optimizando procesos de reporting financiero y operacional
  \item Desarrollé scripts en Visual Basic 6 y SQL para automatizar validaciones de integridad de datos, reduciendo en 50\% el tiempo de cierre mensual contable
  \item Elaboré documentación técnica detallada (diagramas de flujo, diccionarios de datos, manuales de usuario) que facilitó la transferencia de conocimiento y onboarding de nuevos analistas
  \item Realicé extracción, limpieza y transformación de datos (data wrangling) para análisis ad-hoc solicitados por gerencias, utilizando SQL avanzado y Excel con Power Query
  \item Capacité a usuarios finales en herramientas analíticas y procedimientos de parametrización de sistemas ERP, mejorando la autonomía operativa de las áreas de negocio
\end{itemize}
\vspace{2pt}
\textbf{OB GROUP PARK} \hfill Santiago, Chile\\
\textit{Analista de Sistemas y TI} \hfill Mayo 2014 – Abr. 2017
\begin{itemize}
  \item Lideré implementación de soluciones tecnológicas para optimizar procesos contables y de gestión en retail, integrando sistemas de punto de venta con plataformas ERP mediante desarrollo de interfaces en SQL y VBA
  \item Diseñé y mantuve bases de datos relacionales para control de inventario y análisis de ventas, aplicando normalización y modelado entidad-relación que mejoró la consistencia de datos en 45\%
  \item Coordiné proyectos de mejora continua con equipos multidisciplinarios, asegurando compatibilidad entre plataformas legacy y nuevas tecnologías cloud, documentando arquitecturas y flujos de integración
\end{itemize}
\vspace{2pt}
\textbf{Hewlett-Packard Company} \hfill Santiago, Chile\\
\textit{Agente de Mesa de Ayuda Nivel 1} \hfill Ago. 2013 – Mayo 2014
\begin{itemize}
  \item Brindé soporte técnico a usuarios finales diagnosticando y resolviendo incidencias de software y hardware, desarrollando capacidades de troubleshooting sistemático y documentación de soluciones en base de conocimiento
\end{itemize}

\section{Proyectos Destacados}

\textbf{Sistema Inteligente de Recomendación de Productos (E-commerce Analytics)} \hfill 2024
\begin{itemize}
  \item Desarrollé aplicación full-stack con React (TypeScript) y backend en Python (FastAPI) que integra Claude AI para generar recomendaciones personalizadas mediante prompt engineering avanzado y análisis de comportamiento de usuario
  \item Implementé base de datos PostgreSQL optimizada con índices GIN para búsquedas semánticas, logrando tiempos de respuesta <200ms en queries analíticas sobre catálogo de +10,000 productos
  \item Diseñé interfaz UI/UX con componentes reutilizables en React y Tailwind CSS, aplicando principios de accesibilidad (WCAG 2.1) y diseño responsivo, validado con usuarios reales que reportaron 85\% de satisfacción
  \item Automaticé pipeline ETL con Python (pandas, NumPy) para limpieza y enriquecimiento de datos de ventas, alimentando modelo de clustering (K-means) que segmentó clientes en 5 grupos para marketing dirigido
\end{itemize}

\section{Habilidades T\'{e}cnicas}
\textbf{Análisis de Datos \& Machine Learning:} Python (pandas, NumPy, scikit-learn, matplotlib), SQL Server (consultas avanzadas, optimización), Azure SQL, Power BI (DAX, modelado), Excel Avanzado (Power Query, VBA) \\
\textbf{Bases de Datos \& Cloud:} SQL Server (diseño, tuning, T-SQL), Oracle, PostgreSQL, MongoDB, Microsoft Azure (Azure SQL Database, Azure Data Factory), ETL \\
\textbf{Desarrollo \& IA:} Python, C\#, Node.js, TypeScript, Prompt Engineering (Claude, Gemini), APIs REST, JSON/XML \\
\textbf{Herramientas \& Metodologías:} Git (GitHub/GitLab), Docker, SSMS, Visual Studio Code, Scrum/Kanban, Documentación Técnica, Análisis Funcional

\section{Formaci\'{o}n Acad\'{e}mica}
\textbf{Business Analytics Essentials, Machine Learning con Python, Data Science}, Universidad de Chile, Escuela de Postgrado \hfill 2023 \\
\textbf{Optimización de Procesos y Administración Avanzada de Bases de Datos SQL Server}, Universidad de Chile \hfill 2023 \\
\textbf{Desarrollo de Aplicaciones Móviles}, Universidad Complutense de Madrid \hfill 2017 \\
\textbf{Analista Programador}, Universidad Tecnológica de Chile INACAP \hfill 2013

\end{document}